% ---------------------------------------------------------------------------
% Consensus and formation control environment
% ---------------------------------------------------------------------------
\begingroup
% Line breaking for paragraphs with long, unbreakable identifiers (local to
% this chapter): allow looser lines instead of overfull ones.
\setlength{\emergencystretch}{3em}\tolerance=1000
\chapter{Consensus and Formation Control}\label{chap:consensus}

\code{env\_lib.consensus\_env.ConsensusEnv} is a networked multi-agent control
benchmark, new in version 1.0. $N$ agents move in a bounded plane and must
either agree on a common position (\emph{consensus} or rendezvous) or assume a
geometric formation up to translation (\emph{formation control}). Every agent
only sees the agents it is connected to in a communication graph, which is
either fixed or rebuilt every step from the agents' distances. The environment
connects multi-agent reinforcement learning with the classical theory of
distributed control: the Laplacian consensus protocol is included as a baseline
controller, the algebraic connectivity of the graph is reported at every step,
and the difficulty can be tuned through the topology, the dynamics (single or
double integrators), noise and the formation shape. Typical research uses are
learning distributed controllers from local observations, studying the effect of
graph connectivity and time-varying graphs on learning, and comparing learned
policies with model-based protocols.

% ---------------------------------------------------------------------------
\section{Model}\label{sec:consensus-model}

Agent $i$ has position $x_i \in \mathbb{R}^2$, velocity $v_i \in \mathbb{R}^2$
and a formation offset $d_i$. The offsets are zero for consensus and the slots
of the chosen formation shape, centred at the origin, for formation control.
The task is solved when the shifted positions $y_i = x_i - d_i$ agree. The
task error
\begin{equation}
  e = \frac{1}{N} \sum_{i=1}^{N} \lVert y_i - \bar y \rVert^2,
  \qquad \bar y = \frac{1}{N} \sum_{j=1}^{N} y_j,
  \label{eq:consensus-error}
\end{equation}
measures the disagreement; the task is solved when $e <$ \code{tolerance}.

\paragraph{Dynamics.} The action $u_i$ of agent $i$ is clipped to
$[-u_{\max}, u_{\max}]^2$ with $u_{\max}$ = \code{max\_control}. With time step
$\Delta t$ = \code{dt}, noise intensity $\sigma$ = \code{noise\_std},
$w_i \sim \mathcal{N}(0, I_2)$ and the arena
$[-a, a]^2$ ($a$ = \code{arena\_size}):
\begin{itemize}
  \item single integrator (\code{dynamics="single"}, the action is a velocity):
    \begin{equation}
      x_i \leftarrow \operatorname{clip}_{[-a, a]}\bigl(x_i + \Delta t\, u_i
        + \sigma \sqrt{\Delta t}\, w_i\bigr);
    \end{equation}
    the reported velocity is the realised one, $v_i = (x_i^{+} - x_i)/\Delta t$;
  \item double integrator (\code{dynamics="double"}, the action is an
    acceleration, damping $c$ = \code{damping}):
    \begin{equation}
      v_i \leftarrow v_i + \Delta t\,(u_i - c\, v_i) + \sigma \sqrt{\Delta t}\, w_i,
      \qquad
      x_i \leftarrow \operatorname{clip}_{[-a, a]}\bigl(x_i + \Delta t\, v_i\bigr)
    \end{equation}
    (semi-implicit Euler); a velocity component is set to zero when the arena
    wall stops the agent.
\end{itemize}

\paragraph{Communication graph.} The agents communicate over an undirected graph
with adjacency matrix $A$ and neighbour sets $\mathcal{N}_i$
(Table~\ref{tab:consensus-topologies}). Static graphs are built once at
construction; the proximity graph is rebuilt after every step. The
environment reports the algebraic connectivity $\lambda_2(L)$, the second
smallest eigenvalue of the graph Laplacian $L = D - A$ ($D$ the diagonal degree
matrix). It is positive exactly when the graph is connected and bounds the
convergence rate of linear consensus.

\begin{table}[htbp]
  \centering\small
  \caption{Communication topologies.}
  \label{tab:consensus-topologies}
  \begin{tabularx}{\textwidth}{@{}l>{\raggedright\arraybackslash}X@{}}
    \toprule
    \code{topology} & Graph \\
    \midrule
    \code{"ring"} & $i$ is connected to $i \pm 1 \bmod N$ \\
    \code{"line"} & $i$ is connected to $i \pm 1$ (a path) \\
    \code{"star"} & agent 0 is the hub connected to all others \\
    \code{"complete"} & all pairs are connected \\
    \code{"erdos\_renyi"} & every pair is connected independently with probability
      \code{edge\_probability}; the graph is resampled (up to 1000 times) until it is
      connected, using a generator seeded with \code{graph\_seed} \\
    \code{"proximity"} & disk graph: $i$ and $j$ are connected when
      $\lVert x_i - x_j \rVert \le$ \code{comm\_radius}; rebuilt after every step and
      may disconnect during an episode \\
    \bottomrule
  \end{tabularx}
\end{table}

\paragraph{Formation shapes.} \code{make\_formation(shape, n\_agents, radius)}
returns the offsets $d_i$, centred at the origin; all shapes have a
characteristic extent of $2R$ with $R$ = \code{formation\_radius}:
\begin{itemize}
  \item \code{"circle"}: regular polygon of radius $R$, agent 0 at the top;
  \item \code{"line"}: evenly spaced on a horizontal segment of length $2R$;
  \item \code{"grid"}: row-major grid with $\lceil\sqrt{N}\,\rceil$ columns whose
        longer side spans $2R$;
  \item \code{"wedge"}: V shape pointing up, agent 0 at the apex and the others
        alternating between the two arms (half opening angle 35 degrees, arm
        length $2R$).
\end{itemize}
A formation that does not fit into the arena raises \code{ValueError}.

\paragraph{Initial state.} Positions are drawn uniformly from
$[-0.8a, 0.8a]^2$ and velocities are zero. For the proximity topology the draw
is conditioned on a connected initial graph (rejection sampling in vectorised
batches); if \code{comm\_radius} is so small that no connected placement is
found, the agents are placed sequentially, each within $0.95$\,\code{comm\_radius}
of a previously placed agent, and a warning is issued once.
\code{reset(options=\{"positions": ..., "velocities": ...\})} starts from a
given state (arrays of shape \code{(N, 2)}; positions are clipped to the
arena).

% ---------------------------------------------------------------------------
\section{Spaces}\label{sec:consensus-spaces}

\paragraph{Action.} \code{Box(-max\_control, max\_control, (N, 2), float32)}:
one velocity command (single integrator) or acceleration (double integrator) per
agent. Any array with $2N$ entries is reshaped to \code{(N, 2)}; values are
clipped, and NaN or infinite values raise \code{ValueError}.

\paragraph{Observation.} An unbounded \code{Box} of shape \code{(N, 6 + 5K)},
\code{float32}, where $K$ = \code{max\_neighbors} is the number of neighbour
slots. Row $i$ has the layout of Table~\ref{tab:consensus-obs}. The slots hold
the neighbours $j \in \mathcal{N}_i$ sorted by physical distance
$\lVert x_j - x_i \rVert$ (ties by index); unused slots are zero with mask 0.
If an agent has more neighbours than slots, the farthest ones are left out of
the observation; they still count in the reward and in the baseline controller.
The property \code{env.observation\_layout} returns the column slices
(\code{"position"}, \code{"velocity"}, \code{"offset"}, \code{"neighbors"}), and
\code{obs[:, layout["neighbors"]].reshape(N, K, 5)} gives the slots, whose
columns are described by \code{ConsensusEnv.neighbor\_slot\_layout}.

\begin{table}[htbp]
  \centering\small
  \caption{Consensus observation (row $i$; slot $s = 0, \dots, K-1$).}
  \label{tab:consensus-obs}
  \begin{tabularx}{\textwidth}{@{}ll>{\raggedright\arraybackslash}X@{}}
    \toprule
    Columns & Layout key & Content \\
    \midrule
    \code{0:2} & \code{position} & $x_i$ \\
    \code{2:4} & \code{velocity} & $v_i$ \\
    \code{4:6} & \code{offset} & $d_i$ (zeros for consensus) \\
    \code{6+5s : 6+5s+2} & \code{neighbors} & slot $s$: $(x_j - d_j) - (x_i - d_i) = y_j - y_i$ \\
    \code{6+5s+2 : 6+5s+4} & \code{neighbors} & slot $s$: $v_j - v_i$ \\
    \code{6+5s+4} & \code{neighbors} & slot $s$: valid mask (1 or 0) \\
    \bottomrule
  \end{tabularx}
\end{table}

$K$ defaults to the maximum degree of a static graph (2 for the default ring, so
the default observation has shape \code{(8, 16)}) and to $\min(6, N - 1)$ for the
proximity graph.

% ---------------------------------------------------------------------------
\section{Reward and episode end}\label{sec:consensus-reward}

Agent $i$ receives the negative mean squared disagreement with its neighbours
and a quadratic control cost,
\begin{equation}
  r_i = -\frac{1}{\max(|\mathcal{N}_i|, 1)} \sum_{j \in \mathcal{N}_i}
        \lVert y_i - y_j \rVert^2 \;-\; c_u \lVert u_i \rVert^2,
  \label{eq:consensus-reward}
\end{equation}
with $c_u$ = \code{control\_cost} and the clipped action $u_i$, evaluated after
the step on the new graph (isolated agents only pay the control cost). The team
reward is $R = \frac{1}{N} \sum_i r_i$, plus \code{success\_bonus} on the step on
which the task error first drops below \code{tolerance}.

\code{terminated} is \code{True} when $e <$ \code{tolerance};
\code{truncated} is \code{True} once \code{max\_steps} steps were taken. Both
can be \code{True} on the last step.

% ---------------------------------------------------------------------------
\section{Parameters}\label{sec:consensus-params}

All constructor arguments are keyword-only.

\begingroup\small
\begin{longtable}{@{}>{\raggedright\arraybackslash}p{0.24\textwidth}>{\raggedright\arraybackslash}p{0.17\textwidth}>{\raggedright\arraybackslash}p{0.53\textwidth}@{}}
  \caption{Constructor parameters of \code{ConsensusEnv}.}\label{tab:consensus-params}\\
  \toprule
  Parameter & Default & Description \\
  \midrule
  \endfirsthead
  \toprule
  Parameter & Default & Description \\
  \midrule
  \endhead
  \bottomrule
  \endfoot
  \code{n\_agents} & \code{8} & Number of agents $N$ ($\ge 2$). \\
  \code{task} & \code{"consensus"} & \code{"consensus"} (rendezvous) or \code{"formation"}. \\
  \code{topology} & \code{"ring"} & \code{"ring"}, \code{"line"}, \code{"star"}, \code{"complete"}, \code{"erdos\_renyi"} or \code{"proximity"}. \\
  \code{dynamics} & \code{"single"} & \code{"single"} (velocity input) or \code{"double"} (acceleration input with damping). \\
  \code{dt} & \code{0.1} & Integration time step ($> 0$). \\
  \code{max\_steps} & \code{200} & Episode length before truncation. \\
  \code{arena\_size} & \code{10.0} & Positions are clipped to $[-a, a]^2$ ($> 0$). \\
  \code{max\_control} & \code{1.0} & Bound of every action component ($> 0$). \\
  \code{control\_cost} & \code{0.01} & Weight $c_u$ of the control penalty ($\ge 0$). \\
  \code{noise\_std} & \code{0.0} & Noise intensity $\sigma$; each step adds noise with standard deviation $\sigma\sqrt{\Delta t}$ to the positions (single) or velocities (double). \\
  \code{tolerance} & \code{0.05} & Success threshold on the task error ($> 0$, squared distance units). \\
  \code{success\_bonus} & \code{10.0} & Team-reward bonus paid once, on the step the task is solved. \\
  \code{formation\_shape} & \code{"circle"} & \code{"circle"}, \code{"line"}, \code{"grid"} or \code{"wedge"} (formation task only). \\
  \code{formation\_radius} & \code{3.0} & Scale $R$ of the formation ($> 0$). \\
  \code{comm\_radius} & \code{4.0} & Communication range of the proximity topology (inclusive, $> 0$). \\
  \code{edge\_probability} & \code{0.3} & Edge probability of the Erd\H{o}s--R\'enyi topology, in $(0, 1]$. \\
  \code{graph\_seed} & \code{None} & Seed of the Erd\H{o}s--R\'enyi graph, sampled once at construction; \code{None} uses fresh entropy, so two instances may differ in graph and observation size. \\
  \code{max\_neighbors} & \code{None} & Number of neighbour slots $K$ ($\ge 1$); defaults to the maximum degree of a static graph and to $\min(6, N-1)$ for the proximity graph. \\
  \code{damping} & \code{0.5} & Linear velocity damping $c$ of the double integrator ($\ge 0$). \\
  \code{render\_mode} & \code{None} & \code{None}, \code{"human"} or \code{"rgb\_array"}. \\
\end{longtable}
\endgroup

% ---------------------------------------------------------------------------
\section{Info dictionary and properties}\label{sec:consensus-info}

\begin{table}[htbp]
  \centering\small
  \caption{Consensus info dictionary (returned by \code{reset()} and \code{step()}).}
  \label{tab:consensus-info}
  \begin{tabularx}{\textwidth}{@{}l>{\raggedright\arraybackslash}X@{}}
    \toprule
    Key & Content \\
    \midrule
    \code{agent\_rewards} & \code{float64} array \code{(N,)} of the rewards $r_i$ (zeros at reset) \\
    \code{error} & task error $e$ \eqref{eq:consensus-error} \\
    \code{algebraic\_connectivity} & $\lambda_2(L)$ of the current graph \\
    \code{adjacency} & boolean adjacency matrix \code{(N, N)} of the current graph \\
    \code{success} & \code{True} once the task has been solved in the episode (sticky) \\
    \code{step} & number of steps taken in the episode \\
    \bottomrule
  \end{tabularx}
\end{table}

The state is available through read-only properties
(Table~\ref{tab:consensus-properties}); \code{obs\_dim} and
\code{max\_neighbors} are plain attributes.

\begin{table}[htbp]
  \centering\small
  \caption{Read-only properties of \code{ConsensusEnv} (arrays are copies).}
  \label{tab:consensus-properties}
  \begin{tabularx}{\textwidth}{@{}l>{\raggedright\arraybackslash}X@{}}
    \toprule
    Property & Content \\
    \midrule
    \code{positions}, \code{velocities} & $x_i$ and $v_i$, shape \code{(N, 2)} \\
    \code{formation\_offsets} & offsets $d_i$, shape \code{(N, 2)} (zeros for consensus) \\
    \code{adjacency}, \code{laplacian} & adjacency matrix $A$ and Laplacian $L = D - A$ of the current graph \\
    \code{algebraic\_connectivity} & $\lambda_2(L)$ of the current graph \\
    \code{task\_error} & task error $e$ \\
    \code{step\_count} & number of steps taken in the episode \\
    \code{observation\_layout} & column slices of an observation row \\
    \bottomrule
  \end{tabularx}
\end{table}

The module
\code{env\_lib.consensus\_env} also exports the helpers used by the
environment, which are useful for analysis and for reference controllers:
\begin{itemize}
  \item \code{make\_topology(topology, n\_agents, *, edge\_probability=0.3, rng=None)}:
        adjacency matrix of a static graph;
  \item \code{proximity\_adjacency(positions, comm\_radius)}: disk-graph adjacency;
  \item \code{make\_formation(shape, n\_agents, radius=1.0)}: formation offsets;
  \item \code{algebraic\_connectivity(adjacency)} and \code{is\_connected(adjacency)};
  \item the tuples \code{TASKS}, \code{TOPOLOGIES}, \code{DYNAMICS} and
        \code{FORMATION\_SHAPES} of valid option values.
\end{itemize}

% ---------------------------------------------------------------------------
\section{Baseline controller}\label{sec:consensus-baseline}

\code{env.laplacian\_policy(gain=1.0, velocity\_gain=None)} returns the
classical distributed consensus and formation feedback
\begin{equation}
  u_i = -k \sum_{j \in \mathcal{N}_i} \bigl((x_i - d_i) - (x_j - d_j)\bigr),
  \qquad\text{that is}\qquad u = -k\, L\,(x - d),
  \label{eq:laplacian-policy}
\end{equation}
with gain $k$, clipped to the action bounds. For the double integrator the
velocity feedback $-k_v v_i$ is added; the default
$k_v = \max(0, \sqrt{k} - c)$ makes the total damping $c + k_v$ equal to
$\sqrt{k}$. The feedback uses the full neighbour sets, even when the
observation is truncated by \code{max\_neighbors}. The discrete-time
single-integrator loop is stable when $k\, \Delta t\, \lambda_{\max}(L) < 2$;
the method warns once otherwise.

Table~\ref{tab:consensus-baseline} reports the success rate of the controller
with gain 1 for all topologies and task/dynamics combinations (default
parameters otherwise; 100 episodes with reset seeds 0--99 and
\code{graph\_seed} equal to the reset seed). The controller solves every
episode on connected static graphs. On the line, the slowest graph, some
episodes run out of time, and on the proximity graph agents can lose contact.

\begin{table}[htbp]
  \centering\small
  \caption{Success rate of \code{laplacian\_policy(gain=1.0)} over 100 episodes,
    with the mean episode length in steps (truncated episodes count 200).}
  \label{tab:consensus-baseline}
  \begin{tabular}{@{}lcccc@{}}
    \toprule
    & \multicolumn{2}{c}{Consensus} & \multicolumn{2}{c}{Formation (circle)} \\
    \cmidrule(lr){2-3}\cmidrule(l){4-5}
    Topology & single & double & single & double \\
    \midrule
    ring        & 1.00 (84)  & 1.00 (81)  & 1.00 (97)  & 1.00 (86)  \\
    line        & 0.82 (160) & 0.94 (154) & 0.53 (183) & 0.72 (174) \\
    star        & 1.00 (73)  & 1.00 (80)  & 1.00 (85)  & 1.00 (85)  \\
    complete    & 1.00 (65)  & 1.00 (81)  & 1.00 (78)  & 1.00 (87)  \\
    Erd\H{o}s--R\'enyi & 1.00 (90) & 1.00 (90) & 1.00 (104) & 1.00 (97) \\
    proximity   & 0.83 (72)  & 0.70 (109) & 0.55 (135) & 0.12 (190) \\
    \bottomrule
  \end{tabular}
\end{table}

% ---------------------------------------------------------------------------
\section{Rendering}\label{sec:consensus-render}

The dashboard ($1000 \times 560$ pixels, Figures~\ref{fig:consensus} and
\ref{fig:formation}) shows the arena on the left: agents as coloured markers
with halos, communication links whose opacity decreases with their length,
fading trails of the recent positions, the formation slots around the current
centroid (hollow markers, formation task only) and the centroid (cross). On the
right, the task error is plotted on a logarithmic scale with the tolerance as a
dashed line over a shaded success band, and below it the algebraic connectivity
$\lambda_2$. The header reports the task, the team size, the step, the error,
$\lambda_2$ and a status badge.

\begin{figure}[htbp]
  \centering
  \includegraphics[width=\figwidth]{consensus.png}
  \caption{Consensus (rendezvous) of 8 single integrators on a ring, solved at
    step 64 by the Laplacian controller. The agents have gathered at the
    centroid; the task error has fallen below the tolerance and $\lambda_2$ of
    the static ring is constant.}
  \label{fig:consensus}
\end{figure}

\begin{figure}[htbp]
  \centering
  \includegraphics[width=\figwidth]{formation.png}
  \caption{Wedge formation of 8 agents on a proximity graph, solved at step 68.
    Hollow markers are the formation slots around the centroid; $\lambda_2$
    changes as links appear and disappear and drops to zero while the graph is
    disconnected.}
  \label{fig:formation}
\end{figure}

% ---------------------------------------------------------------------------
\section{Registered ids}\label{sec:consensus-ids}

\begin{table}[!htbp]
  \centering\small
  \caption{Registered ids of the consensus environment.}
  \label{tab:consensus-ids}
  \begin{tabular}{@{}lll@{}}
    \toprule
    Id & Class & Configuration \\
    \midrule
    \envid{Consensus-v0} & \code{ConsensusEnv} & defaults (rendezvous, 8 agents on a ring) \\
    \envid{Formation-v0} & \code{ConsensusEnv} & \code{task="formation"} (circle of radius 3) \\
    \bottomrule
  \end{tabular}
\end{table}

Table~\ref{tab:consensus-ids} lists the two ids. Keyword arguments of
\code{env\_lib.make} override the registered configuration, for example \code{env\_lib.make("Formation-v0", formation\_shape="wedge",
topology="proximity")}.

% ---------------------------------------------------------------------------
\section{Usage}\label{sec:consensus-usage}

\begin{pycode}
import env_lib

env = env_lib.make("Formation-v0", formation_shape="wedge",
                   topology="proximity", dynamics="double")
core = env.unwrapped
obs, info = env.reset(seed=0)                 # obs: (8, 36), K = 6 slots
layout = core.observation_layout
n, k = core.n_agents, core.max_neighbors
slots = obs[:, layout["neighbors"]].reshape(n, k, 5)   # neighbour slots
episode_return = 0.0
while True:
    action = core.laplacian_policy(gain=1.0)  # u = -gain L (x - d) - k_v v
    obs, reward, terminated, truncated, info = env.step(action)
    episode_return += reward
    if terminated or truncated:
        break
print(info["step"], info["error"], info["success"], episode_return)
env.close()
\end{pycode}

The example script compares random actions with the Laplacian controller and
records the controller:

\begin{shellcode}
python examples/environments/consensus_demo.py
python examples/environments/consensus_demo.py --task formation --save renders/formation.gif
python examples/environments/consensus_demo.py --task formation --shape wedge \
    --topology proximity --dynamics double --save renders/wedge.gif
\end{shellcode}

% ---------------------------------------------------------------------------
\section{Native vector environment}\label{sec:consensus-vector}

\code{env\_lib.make\_vec("Consensus-v0", B)} and \code{make\_vec("Formation-v0", B)}
return a \code{ConsensusVectorEnv}, which simulates $B$ copies with the same
array kernels as the single environment. It takes the constructor arguments of
\code{ConsensusEnv} plus \code{num\_envs} and \code{autoreset\_mode}, supports
every task, topology (proximity graphs are rebuilt per copy and step),
dynamics and formation shape, and returns observations of shape
$(B, N, d)$, team rewards of shape $(B,)$ and batched info arrays with a
leading $B$ axis. Each copy reproduces a single environment started from the
same state bit for bit:

\begin{pycode}
import numpy as np
import env_lib

envs = env_lib.make_vec("Formation-v0", num_envs=4)
obs, _ = envs.reset(seed=0)                     # (4, 8, 16)
single = env_lib.make("Formation-v0").unwrapped
first, _ = single.reset(options={"positions": envs.unwrapped.positions[0]})
assert np.array_equal(first, obs[0])

action = envs.unwrapped.laplacian_policy()      # batched baseline, (4, 8, 2)
obs, reward, terminated, truncated, info = envs.step(action)
assert np.array_equal(single.step(action[0])[0], obs[0])
\end{pycode}

\code{reset(options=...)} accepts \code{"positions"} and \code{"velocities"}
either for all copies (shape $(N, 2)$) or per copy (shape $(B, N, 2)$), and
combines with \code{"reset\_mask"}. Erdos-Renyi graphs are sampled once per
copy from \code{numpy.random.default\_rng(graph\_seed)}; copy 0 has the graph of
a single environment with the same \code{graph\_seed}. The default
\code{max\_neighbors} is then the maximum degree over all copies, so pass
\code{max\_neighbors} explicitly when observations must have the same width as
a single environment's. \code{render()} draws copy 0. Chapter~\ref{chap:workflow}
covers autoreset modes, seeding and throughput.

\section{Reproducibility and performance}\label{sec:consensus-perf}

The initial positions and the process noise are drawn from
\code{self.np\_random}, seeded by \code{reset(seed=...)}. The Erd\H{o}s--R\'enyi
graph is sampled once at construction from \code{graph\_seed}; pass a seed to
obtain the same graph (and observation size) in every instance, for example in
vectorised training environments.

Measured on the development machine without rendering, a step takes about
0.1\,ms with the default 8 agents and 0.18\,ms for a grid formation of 32
agents on a proximity graph. Rendering an \code{"rgb\_array"} frame takes
about 16--21\,ms.

\endgroup
